\documentclass[11pt]{article}

\usepackage[T1]{fontenc}
\usepackage[utf8]{inputenc}
\usepackage{lmodern}
\usepackage{microtype}
\usepackage[margin=1.05in]{geometry}
\usepackage{amsmath,amssymb,amsthm,mathtools}
\usepackage{booktabs}
\usepackage{array}
\usepackage{float}
\usepackage{enumitem}
\usepackage{xcolor}
\usepackage[numbers,sort&compress]{natbib}
\usepackage[colorlinks=true,linkcolor=blue!55!black,citecolor=blue!55!black,urlcolor=blue!55!black]{hyperref}
\hypersetup{
  pdftitle={Sharp Error Amplification in Recursively Projected Multilinear Computations},
  pdfauthor={Anonymous manuscript}
}

\newtheorem{theorem}{Theorem}[section]
\newtheorem{lemma}[theorem]{Lemma}
\newtheorem{proposition}[theorem]{Proposition}
\newtheorem{corollary}[theorem]{Corollary}
\theoremstyle{definition}
\newtheorem{definition}[theorem]{Definition}
\newtheorem{remark}[theorem]{Remark}

\newcommand{\R}{\mathbb{R}}
\newcommand{\C}{\mathbb{C}}
\newcommand{\K}{\mathbb{K}}
\newcommand{\Ran}{\operatorname{Ran}}
\newcommand{\op}{\mathrm{op}}
\newcommand{\amb}{\mathrm{amb}}
\newcommand{\fin}{\mathrm{fin}}
\newcommand{\proj}{\mathrm{proj}}
\newcommand{\norm}[1]{\left\lVert #1\right\rVert}
\newcommand{\ip}[2]{\left\langle #1,#2\right\rangle}

\setlist[itemize]{leftmargin=1.5em,itemsep=2pt,topsep=4pt}
\setlist[enumerate]{leftmargin=1.7em,itemsep=2pt,topsep=4pt}
\setlength{\parskip}{0.35em}
\setlength{\parindent}{0pt}
\emergencystretch=2em

\title{\textbf{Sharp Error Amplification in Recursively Projected Multilinear Computations}}
\author{Anonymous manuscript}
\date{Draft of \today}

\begin{document}
\maketitle

\begin{abstract}
We study finite-dimensional computations obtained by composing bounded
multilinear maps along a rooted tree and applying an orthogonal projector after
every internal operation.  If each law has operator norm at most $M$ and its
normal output on already-projected inputs is at most $\rho$, then a tree with
$k$ internal vertices has projected root error at most
$(k-1)\rho M^{k-1}$ times the product of the leaf norms.  The missing root
contribution is structural: the final projection removes the root's own normal
defect.  The coefficient is zero for $k=1$ and the normalized constant is one
for $k=2$.  For every ordered tree with exactly three internal vertices,
arbitrary finite arities, independently selectable node laws, finite-dimensional
real or complex Hilbert spaces, and orthogonal projectors, we determine the best
dimension- and rank-independent constant:
\[
 W_3(\eta)=
 \begin{cases}
 \sqrt{4-3\eta^2},&0<\eta\le \sqrt{2/3},\\[1mm]
 \dfrac{2}{\sqrt3\,\eta},&\sqrt{2/3}\le\eta\le1,
 \end{cases}
 \qquad \eta=\rho/M.
\]
The upper bound follows from orthogonal error coupling and a two-dimensional
sum-of-squares identity; matching real two-dimensional constructions attain it.
Thus local error accumulation is first-order sharp as $\eta\downarrow0$, but
the finite-defect extremal geometry is strictly smaller than the universal
coefficient two.  Exact constants at fixed $\eta$ for $k\ge4$ remain open.
All novelty comparisons in this draft are provisional pending theorem-level
bibliographic adjudication and independent review.
\end{abstract}

\section{Introduction}

Projection is a standard way to keep an evolving or recursively constructed
object inside a prescribed reduced space.  Its local effect is transparent:
an orthogonal projector is contractive, and the discarded component can be
measured.  Its global effect is subtler.  Once a discarded component enters a
multilinear composition, subsequent maps can redirect it into the retained
space.  The final retained error is therefore neither the sum of local
discarded norms nor the defect of the root alone.

This paper isolates the corresponding extremal problem.  A finite ordered
rooted tree specifies a composition of multilinear maps.  Every internal value
is projected before being passed upward.  We assume a uniform operator-norm
bound $M$ and a uniform closure-defect bound $\rho$, measured only on inputs
which are already in the appropriate projected ranges.  The question is:
\emph{how large can the retained root error be, relative to the local budget
$\rho$?}

Our first result is the universal estimate
\begin{equation}\label{eq:intro-universal}
 E_T^P\le (k-1)\rho M^{k-1}L_T,
\end{equation}
where $k$ is the number of internal vertices and $L_T$ is the product of leaf
norms.  The main result determines the exact normalized constant for $k=3$ as
the piecewise function $W_3$ displayed in the abstract.  The transition at
$\eta_c=\sqrt{2/3}$ separates a budget-limited regime, in which an extremizer
uses the full defect allowance, from a geometry-limited regime, in which the
worst absolute error has already saturated.

\paragraph{Position relative to neighboring projection theories.}
Four nearby literatures make the scope easier to see.

\begin{itemize}
\item \textbf{HSVD and hierarchical tensor truncation.}  Hierarchical SVD,
hierarchical Tucker, and tensor-train truncation approximate a fixed tensor by
a structured low-rank tensor and compare the truncation with best approximation
in a tensor norm; representative quasi-optimality results are developed in
\cite{Grasedyck2010,Oseledets2011,HackbuschKuhn2009}.  Here the object is instead
a composition of independently bounded multilinear maps, the denominator is a
per-node closure budget, and the comparison is with the unprojected evaluation.
Consequently the familiar $\sqrt{d}$-type truncation factors and our
$(k-1)$/$W_3(\eta)$ amplification constants answer different questions.

\item \textbf{Alternating projections.}  Classical alternating-projection
theory studies products of projectors and convergence to an intersection;
angles between subspaces control sharp error or rate bounds in important cases
\cite{KayalarWeinert1988,DeutschHundal1997}.  Our projectors do not alternate on
one state.  They are interleaved with nonlinear multilinear laws on a tree, and
the propagated component may re-enter a retained range.  Orthogonality remains
decisive, but through Pythagorean coupling at a node rather than through a
Friedrichs angle between target subspaces.

\item \textbf{Tensor-network perturbation and conditioning.}  Sitewise
perturbations of tensor networks already admit worst-case bounds, condition
numbers, and tight witness constructions \cite{ZhangSolomonik2020}; finite-
precision analyses of tree tensor-network operations likewise propagate local
roundoff and truncation errors through a contraction tree
\cite{BaboulinEtAl2025}.  These are substantial precedents for sharp local-to-
global amplification.  Their perturbation variables are tensor sites or
arithmetic operations, whereas our admissibility parameter controls the normal
part of a node law only on already-projected inputs.  They do not determine the
fixed-$\eta$ supremum for the projected-closure error in
\eqref{eq:projected-error}.

\item \textbf{DLRA and TDVP.}  Dynamical low-rank approximation and
time-dependent variational principles project a vector field onto a tangent
space of a fixed-rank manifold and analyze time-integration or variational error
\cite{KochLubich2007,LubichOseledetsVandereycken2015,HaegemanEtAl2016}.  The
present problem is static and finite, has no time step or tangent bundle, and
allows a different multilinear law and projector at every vertex.  It can
serve as a local algebraic model of recursive projection, but it does not by
itself establish a DLRA or TDVP error theorem.
\end{itemize}

\begin{table}[H]
\centering
\small
\caption{Theorem-level positioning from a scoped full-text audit.  Failure to
identify an equivalent theorem in that corpus is not a priority claim.}
\label{tab:positioning}
\begin{tabular}{>{\raggedright\arraybackslash}p{0.25\textwidth}
                >{\raggedright\arraybackslash}p{0.29\textwidth}
                >{\raggedright\arraybackslash}p{0.36\textwidth}}
\toprule
Result in this paper & Closest precedent & Difference requiring adjudication\\
\midrule
Universal projected coefficient $k-1$ & layered/telescoping error bounds;
Zhang--Solomonik sitewise perturbation; TTN finite-precision propagation &
recursive laws with a projected-input closure budget and exact removal of the
root defect; retained root error rather than contracted-tensor perturbation\\
$C_2^P(\eta)=1$ & equality in elementary operator-norm chains & exact class and
closure convention may be unstated folklore\\
$C_3^{P,\fin}(\eta)=W_3(\eta)$ & Zhang--Solomonik tight worst-case site
perturbation; exact alternating-projection constants & nonlinear fixed-defect-
ratio supremum, piecewise phase transition, and attaining projected-law
witness\\
$k=3$ universality over arities and skeletons & tree-based tensor formats &
classification by internal-child skeleton after freezing leaf slots\\
\bottomrule
\end{tabular}
\end{table}

Table~\ref{tab:positioning} is deliberately conservative.  A scoped full-text
audit did not identify an equivalent to the exact $W_3$ statement, but
bibliographic novelty remains unestablished pending independent review.  In
particular, the paper does not claim to introduce sharp tensor-network
stability or local-to-global amplification in general.  Its mathematical
center is the exact fixed-$\eta$ projected-closure constant.  The arity and
skeleton universality below is a structural consequence of the two sharp
three-node cases, not a separate broad novelty claim.

\paragraph{Contributions.}
\begin{enumerate}
\item We prove the structural bound \eqref{eq:intro-universal} directly by
subtree induction and argument-wise telescoping.
\item We identify the exact constants at $k=1,2$ in the declared class.
\item For $k=3$ we derive a contraction Gram lemma, reduce both possible tree
skeletons to the same two-variable extremal problem, and solve it by a
sum-of-squares identity.
\item We give explicit real two-dimensional witnesses and extend the result to
every finite arity profile by freezing additional leaf slots.
\item We state the fixed-$\eta$ problem for $k\ge4$ as open, rather than
extrapolating from the universal upper bound or from small cases.
\end{enumerate}

\section{Projected multilinear trees}

Let $T$ be a finite ordered rooted tree.  Its internal vertices are
$V_{\mathrm{int}}(T)$, its leaves are $L(T)$, and
$k=k(T)=|V_{\mathrm{int}}(T)|$.  Every vertex $w$ carries a finite-dimensional
Hilbert space $H_w$ over $\K\in\{\R,\C\}$.  An internal vertex $v$ has ordered
children $c_1(v),\ldots,c_{m_v}(v)$, with $m_v\ge2$, and a multilinear law
\[
 \mu_v:H_{c_1(v)}\times\cdots\times H_{c_{m_v}(v)}\longrightarrow H_v.
\]
Every internal $v$ also carries an orthogonal projector
$P_v=P_v^*=P_v^2$.  For a leaf $\ell$ set $\widehat P_\ell=I_{H_\ell}$, and
for an internal vertex set $\widehat P_v=P_v$.

Given leaf data $z_\ell\in H_\ell$, define the ambient evaluation recursively
by $F_\ell=z_\ell$ and
\[
 F_v=\mu_v(F_{c_1(v)},\ldots,F_{c_{m_v}(v)}).
\]
The recursively projected evaluation is $R_\ell=z_\ell$ and
\[
 R_v=P_v\mu_v(R_{c_1(v)},\ldots,R_{c_{m_v}(v)}).
\]
For the root $r$, the quantity of interest is
\begin{equation}\label{eq:projected-error}
 E_T^P=\norm{P_rF_r-R_r}.
\end{equation}
This differs from the ambient discrepancy $\norm{F_r-R_r}$: it retains only
the component visible in the reduced root space.

\begin{definition}[Projected-input closure defect]
For an internal vertex $v$, set
\begin{equation}\label{eq:defect}
 \rho_v^{\proj}=
 \left\|(I-P_v)\mu_v
 \bigl(\widehat P_{c_1(v)}\,\cdot,\ldots,
       \widehat P_{c_{m_v}(v)}\,\cdot\bigr)\right\|_{\op}.
\end{equation}
\end{definition}

The restriction in \eqref{eq:defect} is load-bearing.  It controls the normal
output when the children have already been projected, exactly the inputs seen
by the recursive evaluation.  It is weaker than the unrestricted ambient
defect $\norm{(I-P_v)\mu_v}_{\op}$, and the two conventions must not be
interchanged.

A realization is $(M,\rho)$-admissible if
$\norm{\mu_v}_{\op}\le M$ and $\rho_v^{\proj}\le\rho$ at every internal
vertex, where $M>0$ and $0<\rho\le M$.  Write
\[
 \eta=\rho/M\in(0,1],\qquad
 L_T=\prod_{\ell\in L(T)}\norm{z_\ell}.
\]
For a fixed tree $T$, define the dimension- and rank-independent constant
over independently selectable node laws by
\begin{equation}\label{eq:constant}
 C_T^{P,\fin}(\eta)=
 \sup\frac{E_T^P}{\rho M^{k-1}L_T},
\end{equation}
where the supremum ranges over all nondegenerate admissible
finite-dimensional realizations with $\rho/M=\eta$.  Equality in
\eqref{eq:constant} is a class statement; it does not say that every fixed
dimension and projector rank attains the supremum.

Scaling all node laws by a common positive factor and each leaf independently
shows that the ratio in \eqref{eq:constant} is invariant.  We may therefore
normalize $M=1$ and all nonzero leaves to unit norm.  A zero leaf makes the
entire multilinear evaluation vanish.

\section{The universal projected bound}

For a vertex $v$, let $T_v$ be its subtree, let $k_v$ be its number of
internal vertices, and let $L_v$ be the product of its leaf norms.

\begin{lemma}[Subtree stability]\label{lem:stability}
For every vertex $v$,
\[
 \norm{F_v}\le M^{k_v}L_v,
 \qquad
 \norm{R_v}\le M^{k_v}L_v.
\]
\end{lemma}
\begin{proof}
The statement is immediate at a leaf.  At an internal vertex it follows by
multilinearity, the operator-norm bound on $\mu_v$, contractivity of $P_v$,
and induction over the child subtrees.
\end{proof}

\begin{lemma}[Ambient subtree discrepancy]\label{lem:ambient-subtree}
For every vertex $v$,
\[
 \norm{F_v-R_v}\le k_v\rho M^{k_v-1}L_v.
\]
\end{lemma}
\begin{proof}
The result is trivial at a leaf.  Let $v$ be internal and abbreviate the child
values by $F_i,R_i$.  Add and subtract $\mu_v(R_1,\ldots,R_{m_v})$:
\[
 F_v-R_v=
 \bigl[\mu_v(F_1,\ldots,F_{m_v})-
       \mu_v(R_1,\ldots,R_{m_v})\bigr]
 +(I-P_v)\mu_v(R_1,\ldots,R_{m_v}).
\]
The last term is bounded by
$\rho M^{k_v-1}L_v$ because every $R_i$ is in the corresponding projected
range.  Telescoping the first bracket one argument at a time gives one term
for each child.  Lemma~\ref{lem:stability} and the inductive hypothesis bound
the $i$th term by $k_{c_i}\rho M^{k_v-1}L_v$.  Since
$k_v=1+\sum_i k_{c_i}$, summing proves the claim.
\end{proof}

\begin{theorem}[Structural $k-1$ bound]\label{thm:kminusone}
Every admissible tree with $k\ge1$ internal vertices satisfies
\begin{equation}\label{eq:kminusone}
 \boxed{E_T^P\le(k-1)\rho M^{k-1}L_T.}
\end{equation}
Consequently $C_T^{P,\fin}(\eta)\le k-1$.
\end{theorem}
\begin{proof}
At the root,
\[
 P_rF_r-R_r=
 P_r\bigl[\mu_r(F_1,\ldots,F_{m_r})-
           \mu_r(R_1,\ldots,R_{m_r})\bigr].
\]
The root's local normal defect is absent because $P_r(I-P_r)=0$.
Telescoping over the root arguments and applying
Lemma~\ref{lem:ambient-subtree} to each child yields
\[
 E_T^P\le\Bigl(\sum_i k_{c_i}\Bigr)
 \rho M^{k-1}L_T=(k-1)\rho M^{k-1}L_T.
\]
\end{proof}

The improvement from $k$ to $k-1$ is therefore not a cancellation inferred
from examples.  It follows from which component the final projector removes.

\section{The cases \texorpdfstring{$k=1$ and $k=2$}{k=1 and k=2}}

\begin{corollary}\label{cor:k1}
For $k=1$, $C_1^{P,\fin}(\eta)=0$.
\end{corollary}
\begin{proof}
There is only the root, and by definition $R_r=P_rF_r$.
\end{proof}

For $k=2$, after freezing any additional unit leaf slots, the internal-child
skeleton is a chain.  In the binary notation, let the first law act on leaves
$a,b$, let the second law act on the first value and a leaf $d$, and put
\[
 D=(I-P_1)\mu_1(a,b),\qquad R_1=P_1\mu_1(a,b).
\]
Then
\begin{equation}\label{eq:k2-identity}
 E_T^P=\norm{P_2\mu_2(D,d)}.
\end{equation}

\begin{theorem}[Exact two-vertex constant]\label{thm:k2}
For the unrestricted finite-dimensional independent-law class,
\[
 \boxed{C_2^{P,\fin}(\eta)=1\qquad(0<\eta\le1).}
\]
In a fixed binary realization, equality in the universal estimate holds if and
only if
\begin{align*}
 \norm{D}&=\rho\norm a\norm b,\\
 \norm{\mu_2(D,d)}&=M\norm D\norm d,\\
 \mu_2(D,d)&\in\Ran(P_2).
\end{align*}
\end{theorem}
\begin{proof}
Equation~\eqref{eq:k2-identity} gives the upper bound and the three equality
conditions by inspecting the chain
\[
 E_T^P\le\norm{\mu_2(D,d)}
 \le M\norm D\norm d
 \le \rho M\norm a\norm b\norm d.
\]
For attainment, normalize $M=1$ and take $H=\R^2$ with basis $e_0,e_1$ and
$P$ the projection onto $\R e_0$.  With unit leaves $e_0$, choose
\[
 \mu_1(x,y)=\ip{x}{e_0}\ip{y}{e_0}
 (\sqrt{1-\eta^2}\,e_0+\eta e_1),
 \qquad
 \mu_2(x,y)=\ip{x}{e_1}\ip{y}{e_0}e_0.
\]
Both laws have norm at most one, the first closure defect is $\eta$, the root
law has zero normal output, and $E_T^P=\eta$.  Restoring scale gives
$E_T^P=\rho M L_T$.  Unit gates embed the same witness into additional leaf
slots.
\end{proof}

\section{Three vertices: orthogonal coupling}

The universal coefficient at $k=3$ is two.  It is not the exact finite-defect
constant.  Orthogonality couples the size of a discarded component to the size
of the retained value from which it came.

At the first non-root internal vertex, write
\[
 D_1=(I-P_1)F_1,\qquad R_1=P_1F_1.
\]
Then $D_1\perp R_1$ and, in the normalized setting,
\begin{equation}\label{eq:pythagorean}
 \norm{D_1}^2+\norm{R_1}^2=\norm{F_1}^2\le1.
\end{equation}

The second ingredient controls the cross term created when a contraction acts
on two orthogonal directions.

\begin{lemma}[Projector-contraction Gram bound]\label{lem:gram}
Let $N:H\to K$ be linear with $\norm N_{\op}\le1$, let $u,v\in H$ be
orthonormal, and let $Q$ be an orthogonal projector on $K$.  If
$s=\norm{QNv}$, then
\begin{equation}\label{eq:gram-bound}
 \left|\ip{Nu}{QNv}\right|\le s\sqrt{1-s^2}.
\end{equation}
\end{lemma}
\begin{proof}
Set $A=N^*QN$.  Then $0\preceq A\preceq I$.  On
$\operatorname{span}\{u,v\}$ write $A_{vv}=s^2$, $A_{uu}=\alpha$, and
$A_{uv}=c$.  Positivity of $A$ gives $|c|^2\le\alpha s^2$, while positivity
of $I-A$ gives $|c|^2\le(1-\alpha)(1-s^2)$.  The largest possible minimum of
these two quantities occurs at $\alpha=1-s^2$, proving
$|c|\le s\sqrt{1-s^2}$.
\end{proof}

\subsection{The chain skeleton}

Normalize $M=L_T=1$ and $\rho=\eta$.  Put
$q=\norm{D_1}\le\eta$ and $p=\norm{R_1}\le\sqrt{1-q^2}$.  Freezing the leaf
inputs at the next vertex reduces its law to a linear contraction $N$.  With
$Q=I-P_2$ and suitable orthonormal $u,v$, put $s=\norm{QNv}\le\eta$.
The discrepancy that can reach the projected root has the form
$qNu+pQNv$.  Lemma~\ref{lem:gram} therefore gives
\begin{equation}\label{eq:chain-objective}
 \norm{qNu+pQNv}^2
 \le q^2+p^2s^2+2qps\sqrt{1-s^2}.
\end{equation}
The right-hand side increases with $p$, so take
$p=\sqrt{1-q^2}$.  Write $q=\sin\alpha$, $s=\sin\beta$ and
$\xi=\tan\alpha$, $\zeta=\tan\beta$.  The squared objective becomes
\begin{equation}\label{eq:f-objective}
 f(\xi,\zeta)=
 \frac{(\xi+\zeta)^2+\xi^2\zeta^2}
 {(1+\xi^2)(1+\zeta^2)}.
\end{equation}
The extremal calculation is encoded in the identity
\begin{equation}\label{eq:sos}
 4(1+\xi^2)(1+\zeta^2)
 -3\bigl[(\xi+\zeta)^2+\xi^2\zeta^2\bigr]
 =(\xi\zeta-2)^2+(\xi-\zeta)^2.
\end{equation}
Hence $f\le4/3$, with equality exactly at
$\xi=\zeta=\sqrt2$, or $q=s=\sqrt{2/3}$.  Under the box constraints
$q,s\le\eta$, the maximum is at $q=s=\eta$ until that point becomes
available and remains $4/3$ thereafter.  Dividing the maximal absolute error
by $\eta$ gives
\begin{equation}\label{eq:w3}
 W_3(\eta)=
 \begin{cases}
 \sqrt{4-3\eta^2},&0<\eta\le\sqrt{2/3},\\[1mm]
 \dfrac{2}{\sqrt3\,\eta},&\sqrt{2/3}\le\eta\le1.
 \end{cases}
\end{equation}

\begin{theorem}[Exact chain constant]\label{thm:chain}
For the ordered three-vertex chain with independently selectable laws,
\[
 C_{3,\mathrm{chain}}^{P,\fin}(\eta)=W_3(\eta).
\]
\end{theorem}
\begin{proof}
The preceding argument gives the upper bound.  For attainment take
$H=\R^2$, $P=e_0e_0^*$, all leaves equal to $e_0$, and
\[
 t=\min\{\eta,\sqrt{2/3}\},\qquad p=c=\sqrt{1-t^2}.
\]
Define
\begin{align*}
 \mu_1(x,y)&=\ip{x}{e_0}\ip{y}{e_0}(te_1+pe_0),\\
 \mu_2(x,y)&=\ip{y}{e_0}N_tx,
\end{align*}
where the orthogonal map $N_t$ satisfies
\[
 N_te_0=te_1-ce_0,\qquad N_te_1=ce_1+te_0.
\]
After two levels the propagated discrepancy is
\[
 v=tN_te_1+pt,e_1=t^2e_0+2tp,e_1,
 \qquad \norm v=t\sqrt{4-3t^2}.
\]
Let $Ah=\ip{v}{h}\norm v^{-1}e_0$ and
$\mu_3(x,y)=\ip{y}{e_0}Ax$.  All three laws are contractions, both active
defects equal $t\le\eta$, and the root defect is zero.  The projected root
error is $\norm v$, which gives \eqref{eq:w3} after division by $\eta$.
\end{proof}

\begin{remark}[Two regimes]
The absolute normalized extremal error is
\[
 G_3(\eta):=\eta W_3(\eta)=
 \begin{cases}
 \eta\sqrt{4-3\eta^2},&\eta\le\sqrt{2/3},\\[1mm]
 2/\sqrt3,&\eta\ge\sqrt{2/3}.
 \end{cases}
\]
Below the transition the witness uses the full defect budget.  Above it, the
effective leakage freezes at $\sqrt{2/3}$; the decrease of $W_3$ results from
normalizing a saturated absolute error by a growing allowed budget.
\end{remark}

\subsection{The branching skeleton}

Now two internal children feed the root.  Write their discarded and retained
norms as $q_i$ and $p_i$, with
$q_i\le\eta$ and $p_i\le\sqrt{1-q_i^2}$.  After scalarizing the projected root
output, maximizing over contractive bilinear forms is dual to a nuclear norm.
The relevant two-by-two coefficient matrix is
\[
 A=\begin{pmatrix}
 q_1q_2&q_1p_2\\
 p_1q_2&0
 \end{pmatrix}.
\]
Since $\norm A_*^2=\norm A_F^2+2|\det A|$, saturation of the $p_i$ bounds gives
\begin{equation}\label{eq:branch-objective}
 \norm A_*^2=
 q_1^2+q_2^2-q_1^2q_2^2
 +2q_1q_2\sqrt{1-q_1^2}\sqrt{1-q_2^2}.
\end{equation}
With $q_1=\sin\alpha$, $q_2=\sin\beta$, expression
\eqref{eq:branch-objective} is exactly $f(\tan\alpha,\tan\beta)$ from
\eqref{eq:f-objective}.  Thus the same sum-of-squares identity controls both
skeletons.

\begin{theorem}[Exact branching constant]\label{thm:branch}
For the ordered three-vertex branching tree with independently selectable
laws,
\[
 C_{3,\mathrm{branch}}^{P,\fin}(\eta)=W_3(\eta).
\]
\end{theorem}
\begin{proof}
Equation~\eqref{eq:branch-objective} and
\eqref{eq:sos} give the upper bound.  For attainment, use two copies of the
first-node output $te_1+pe_0$ from Theorem~\ref{thm:chain}, with
$t=\min\{\eta,\sqrt{2/3}\}$.  At the root choose a norm-one bilinear scalar
form given by a polar factor of the coefficient matrix $A$, and direct that
scalar into $\Ran(P_r)$.  Operator/nuclear-norm duality then attains
$\norm A_*=t\sqrt{4-3t^2}$, while the root closure defect is zero.  This is a
real two-dimensional rank-one-projector realization and yields $W_3(\eta)$.
\end{proof}

\section{Structural universality at \texorpdfstring{$k=3$}{k=3}}

Deleting the leaves from a rooted tree with exactly three internal vertices
leaves only two possible internal-child skeletons: a chain, or two siblings
feeding the root.  Higher arity adds only leaf-occupied slots.

Fixing a unit leaf in such a slot converts a multilinear law into an effective
law on its internal-child slots.  This operation does not increase the operator
norm.  Because leaves carry the identity extended projector, it also does not
increase the projected-input closure defect.  The effective law is linear at a
vertex with one internal child and bilinear at the root of the branching
skeleton.  Conversely, the two-dimensional witnesses above embed in every
finite arity profile by inserting unit $e_0$ gates in the extra leaf slots.

\begin{theorem}[Universal exact constant for three vertices]\label{thm:k3-universal}
Let $T$ be any finite ordered rooted tree with exactly three internal vertices
and internal arities at least two.  Then every admissible finite-dimensional
real or complex realization satisfies
\[
 E_T^P\le W_3(\eta)\rho M^2L_T.
\]
Moreover, over the unrestricted finite-dimensional class with independently
selectable node laws,
\begin{equation}\label{eq:k3-universal}
 \boxed{C_T^{P,\fin}(\eta)=W_3(\eta)\qquad(0<\eta\le1).}
\end{equation}
The equality is attained over $\R$ in dimension two with rank-one projectors.
\end{theorem}
\begin{proof}
Freeze unit leaf slots to reduce to the chain or branching analysis, giving the
upper bound without increasing either admissibility budget.  Embed the matching
two-dimensional witness into the original arity profile with unit gates to
obtain the lower bound.
\end{proof}

\begin{remark}[Quantifiers]
Theorem~\ref{thm:k3-universal} determines the best constant uniform in finite
dimension, rank, arity, and skeleton.  It does not assert equality in each
fixed realization.  For example, if every $P_v=I$, then the projected error is
identically zero.
\end{remark}

For $0<\eta<\sqrt{2/3}$,
\begin{equation}\label{eq:expansion}
 W_3(\eta)=2-\frac34\eta^2-\frac{9}{64}\eta^4+O(\eta^6).
\end{equation}
Thus the universal coefficient two is asymptotically correct to first order,
while orthogonality forces a strict second-order improvement at every positive
defect ratio.

\section{What remains open}

Theorem~\ref{thm:kminusone} applies to every finite tree, but an upper bound is
not an exact constant.  For a fixed tree with $k\ge4$ internal vertices, the
central unresolved finite-defect problem is
\[
 \boxed{C_T^{P,\fin}(\eta)=?\qquad(0<\eta\le1).}
\]
The $k=3$ proof does not extend by independent repetition: a larger tree couples
several retained and discarded components, and the two-variable Gram geometry
is replaced by a higher-dimensional compatibility problem.  The following
questions are therefore left open.

\begin{enumerate}
\item Determine $C_T^{P,\fin}(\eta)$ at fixed $\eta$ for $k=4$, separately for
the distinct internal-child skeletons.
\item Decide whether the leading finite-defect correction to $k-1$ is always
quadratic in $\eta$ and whether its coefficient is a combinatorial invariant of
$T$.
\item Determine how exact constants change under law sharing, fixed dimensions
or projector ranks, and oblique rather than orthogonal projectors.
\end{enumerate}

These are mathematical questions about the declared extremal class.  Numerical
search may suggest witnesses, but cannot establish a universal constant.

\section{Limitations and novelty boundary}

The results are finite-dimensional and deterministic.  They require orthogonal
projectors, independently selectable node laws for the unrestricted sharpness
statements, and the projected-input closure convention
\eqref{eq:defect}.  They do not imply continuum limits, time-integration error
bounds or tensor-rank quasi-optimality.

The bibliography records a scoped full-text comparison with HSVD/HT,
tensor-network perturbation, alternating projections, and DLRA/TDVP.  No
equivalent to \eqref{eq:k3-universal} was identified in that corpus, but this
negative search result is not a proof of novelty.  The citation neighborhood
of sharp tensor-network conditioning, the original Kayalar--Weinert paper, and
the approximate-multiplicativity literature still require independent human
review.  The appropriate claim is therefore \emph{we determine the exact
constant in the declared model}, not a priority claim.

\section{Conclusion}

Recursive projection through a multilinear tree has a simple universal budget
and a nontrivial exact geometry.  The retained root error receives no local
contribution from the root itself, which yields the coefficient $k-1$.  At
three internal vertices, however, orthogonality prevents simultaneous
saturation of the two propagated sources.  Both possible skeletons reduce to
the same sum-of-squares problem, giving the exact constant $W_3(\eta)$ and a
phase transition at $\eta=\sqrt{2/3}$.  Determining the corresponding
fixed-defect constants for $k\ge4$ is the next mathematical problem.

\bibliographystyle{plainnat}
\bibliography{references}

\end{document}
